\section{Stable torus quotients and invariant sections}\label{sec:quotient}

Let $Z$ be smooth and projective, let $L=-K_Z$ have a smooth semipositive Hermitian metric, and let $T$ be an algebraic torus acting on $Z$. We construct a projective quotient with a klt stabilizer boundary and a bounded semipositive anti-log-canonical metric. Rational connectedness is needed only in the final application of Theorem~\ref{thm:klt}, where it gives Euler characteristic one on a resolution of the quotient.

Replace $T$ by its effective image in $\Aut(Z)$. The kernel acts trivially on the tangent bundle and hence on the natural linearization of $L$, so this does not change invariants. For the quotient construction assume that $t=\dim T>0$; the trivial-action case of Theorem~\ref{thm:equivariant} will be treated at the end. Average the weights of the original metric over the compact form $K\simeq(S^1)^t$ of $T$. This gives a smooth $K$-invariant semipositive metric $h$ on the same naturally linearized line bundle.

\subsection{The weight convention and an ample perturbation}\label{subsec:weights}

For a $T$-linearized line bundle $C$, a point $x\in Z$, and an integral one-parameter subgroup $\lambda:\CC^*\to T$, define
\begin{equation}\label{eq:weight}
w_C(x,\lambda)=\text{the weight on }C_{x_0},
\qquad x_0=\lim_{\tau\to0}\lambda(\tau)x.
\end{equation}
Projectivity gives the limit. Our convention is that a fiber vector of weight $w$ transforms by multiplication by $\tau^w$. For a very ample linearized bundle $C$, this is the maximum of the pairings with $\lambda$ of the weights of eigensections that do not vanish at $x$.

To verify the sign, let $s$ be an eigensection of weight $\alpha$ and $l$ a nonzero lift at $x$. Then the scalar of $s$ relative to $\lambda(\tau)l$ equals $\tau^{-\alpha}$ times its scalar at $x$. The smallest order among these scalars is $-\max\alpha$, and a basepoint-free collection has minimum order zero relative to a nonvanishing frame at the limit. Thus the transformed lift has order $\max\alpha$, agreeing with \eqref{eq:weight}. In this convention the Hilbert--Mumford condition is $w_C(x,\lambda)\ge0$ for semistability and strict positivity for stability, for every nonzero one-parameter subgroup; see \cite{MFK1994} and \cite[\S1.1.5]{DolgachevHu1998}.

Torus linearization and equivariant embeddings \cite{Sumihiro1974}, in the precise forms \cite[Proposition~2.4 and Corollary~2.6]{KKLV1989}, allow an equivariantly very ample line bundle $J$. Replace it by a sufficiently high power so that $J+L$ is also very ample, using the natural linearization on $L$. For each of the two very ample bundles, the nonvanishing eigensections at a point form a subset of a fixed finite list. Consequently their weight functions, and the difference
\[
w_L=w_{J+L}-w_J,
\]
have only finitely many patterns. Each pattern is homogeneous and piecewise linear on a finite rational polyhedral fan in $\Lie(T)_{\RR}$, interpreted as the real cocharacter space.

\begin{lemma}\label{lem:genericweight}
On a dense Zariski open subset of $Z$, the generic weight pattern of $L$ is strictly positive on every nonzero real cocharacter.
\end{lemma}
\begin{proof}
First fix a nonzero integral cocharacter $\lambda$. The Bia\l{}ynicki-Birula decomposition \cite{BB1973} has an open attracting stratum; the smooth attracting-piece statement and its tangent-weight description also follow from \cite[Theorem~1.5 and Corollary~7.3]{JS2018}. At its limiting fixed component, all tangent weights are nonnegative, and some are positive. Otherwise the attracting stratum would have no positive-dimensional directions transverse to a fixed component and openness would make $\lambda$ act trivially. This contradicts effectiveness. The weight of $L=\det T_Z$ is the sum of those tangent weights, so it is strictly positive on this open stratum.

The generic pattern for $J$ and $J+L$ occurs on one dense open subset, obtained by requiring every nonzero eigensection in chosen bases to be nonvanishing. This open meets the open attracting stratum for each fixed $\lambda$. Its weight function is therefore positive for every nonzero integral cocharacter. The finitely many rational linear pieces then make it positive on every nonzero real cocharacter: a zero or negative locus on one rational cone would contain a nonzero rational point. On a unit sphere this continuous homogeneous function has a positive minimum.
\end{proof}

\begin{proposition}\label{prop:stableopen}
There are a positive rational number $\delta$ and a rational character shift of arbitrarily small size for the ample rational linearized bundle $L+\delta J$ such that its stable locus $U$ is nonempty and equals its semistable locus. Every $x\in U$ satisfies
\begin{equation}\label{eq:weakweights}
w_L(x,\lambda)\ge0
\quad\text{for all integral one-parameter subgroups }\lambda.
\end{equation}
The shift is used only to define $U$; the linearization on $L$ remains natural.
\end{proposition}
\begin{proof}
The bundle $L$ is nef, so $L+\delta J$ is ample for every $\delta>0$. Subdivide the finitely many weight fans to one finite pointed rational fan. On each cone every weight pattern of $J$ and $L$ is linear.

For each pattern failing \eqref{eq:weakweights}, choose an integral ray on which its weight is negative. There are only finitely many choices. For sufficiently small $\delta>0$ and sufficiently small rational character shift, all these chosen weights remain negative. Such a pattern cannot occur at a semistable point for the perturbed linearization. On the other hand, Lemma~\ref{lem:genericweight} and its uniform positive lower bound on the unit sphere show that the generic pattern stays strictly positive under these perturbations. Thus the stable locus is nonempty.

After fixing such a rational $\delta$, choose the rational character off the finitely many hyperplanes on which a pattern has zero weight on one of the rays of the common fan. For any semistable point, weights on those rays are now strictly positive. Linearity on pointed cones makes them strictly positive on every nonzero cocharacter. The point is stable by Hilbert--Mumford. Finally clear denominators to obtain an ordinary ample linearized line bundle. This leaves the stated loci unchanged.
\end{proof}

By projective GIT the quotient
\begin{equation}\label{eq:quotient}
q:U\longrightarrow W=U/T
\end{equation}
is a geometric quotient and $W$ is projective. The action has finite stabilizers on $U$. Equality of stable and semistable loci is important here: the quotient is the complete projective GIT quotient, not merely an open subset of one.

\subsection{The stabilizer boundary and the anticanonical identity}\label{subsec:slices}

We describe the singularities of $W$ and identify the divisor to which Theorem~\ref{thm:klt} will apply.

\begin{proposition}\label{prop:boundary}
The variety $W$ is normal and $\QQ$-factorial. There is an effective rational divisor
\[
D=\sum_B\left(1-\frac1{e_B}\right)B
\]
on $W$, where $e_B$ is the generic ramification index of a finite stabilizer slice over the prime $B$, such that $(W,D)$ is klt. If $P=-(K_W+D)$, then
\begin{equation}\label{eq:canonicalquotient}
L|_U=q^*P
\end{equation}
as equivariant $\QQ$-line bundles, using the natural linearization on $L$ and the trivial action on the base.
\end{proposition}

\begin{proof}
Apply Luna's slice theorem \cite{Luna1973} in an invariant affine neighborhood coming from the ample GIT linearization. Near an orbit with finite stabilizer $G$, there is a smooth slice $S_{\mathrm{sl}}$ and, after an \emph{\'etale} base localization, the local model
\[
T\times^G S_{\mathrm{sl}}\longrightarrow S_{\mathrm{sl}}/G.
\]
The slice is smooth because $U$ is smooth and the orbit has constant dimension. Normality follows from invariants of a normal ring. Finite quotients of smooth varieties are locally $\QQ$-factorial: a Weil divisor pulls back to a Cartier divisor locally on the smooth finite cover, and the product of its translates gives a local equation for a positive multiple downstairs. The same assertion descends through the \emph{\'etale} localizations. Thus $W$ is $\QQ$-factorial.

Discard an ineffective kernel of the finite slice action when computing ramification. At a general point over a prime $B$, the finite map has ramification order $e_B$, and the ramification divisor has coefficient $e_B-1$. Thus its canonical divisor formula is
\[
K_{S_{\mathrm{sl}}}=p^*(K_W+D)
\]
in codimension one. The numbers $e_B$ are intrinsic: they are the multiplicities along components of the quotient pullback of $B$. The formulas therefore agree on overlapping charts and define one global effective rational divisor $D$. Only finitely many primes are ramified.

The usual finite-morphism discrepancy formula shows that this pair is klt. More explicitly, extend any divisorial valuation over $W$ to the function field of the slice. If the valuation has ramification index $e$, its log discrepancy upstairs is $e$ times its log discrepancy for $(W,D)$ downstairs. The upstairs variety is smooth with zero boundary, hence every log discrepancy is positive. The downstairs log discrepancies are therefore positive as well. This is the finite-morphism discrepancy criterion of \cite[Proposition~5.20]{KM1998}; the exact finite-group calculation with its ramification boundary is also given in \cite[Proposition~3.5 and equation~(3.4)]{BGLM2024}.

It remains to keep track of the linearization in \eqref{eq:canonicalquotient}. On $U$ the infinitesimal action gives an injective subbundle
\[
\Lie(T)\otimes\OO_U\ \longrightarrow\ T_U.
\]
Its rank is $\dim T$ because all stabilizers are finite. Choose a dense smooth open $W^\circ\subset W$ away from the branch boundary over which $q$ is a submersion, and put $U^\circ=q^{-1}(W^\circ)$. There the tangent sequence is
\[
0\longrightarrow\Lie(T)\otimes\OO_{U^\circ}
\longrightarrow T_{U^\circ}
\longrightarrow q^*T_{W^\circ}\longrightarrow0.
\]
The adjoint representation of a torus is trivial, so a nonzero element of $\det\Lie(T)$ gives an equivariant trivialization of the orbit determinant. Taking determinants of this sequence therefore identifies $L|_{U^\circ}$ with $q^*(-K_{W^\circ})$ equivariantly. On each slice the orders of this rational identification are exactly those in the ramification calculation above, which introduces the boundary $D$. No character twist occurs.

The quotient has constant fiber dimension $\dim T$. A prime divisor in $U$ cannot map into codimension at least two in $W$: the dimension bound for such a preimage would be at most $\dim U-2$. Thus the preceding slice and codimension-one calculation accounts for every prime divisor upstairs. It gives the equivariant rational-line-bundle identity \eqref{eq:canonicalquotient}. Equivalently, choose finitely many slice charts covering the projective quotient. Their finite stabilizer groups have bounded exponent, and a common multiple also clears the Cartier indices. At that multiple we obtain an actual isomorphism of line bundles whose action on the pullback from $W$ is trivial along each fiber. This is an equivariant rational-line-bundle identity, not only a numerical equality.
\end{proof}

\subsection{A uniform bound on translated lifts}\label{subsec:orbitbound}

The quotient is now algebraically suitable. To apply Theorem~\ref{thm:klt} and later extend its section, we need a locally bounded quotient metric. We prove the required analytic estimate instead of inferring it from the existence of the GIT quotient.

\begin{lemma}[Orbit bound]\label{lem:orbitbound}
Let $x_0\in U$ and let $l(x)$ be a nonvanishing algebraic local frame of $L$ near $x_0$. Then
\begin{equation}\label{eq:orbitbound}
\sup_{t\in T}|t\cdot l(x)|_h
\end{equation}
is bounded above on a neighborhood of $x_0$. The same holds for holomorphic local frames and positive tensor powers of $L$.
\end{lemma}

\begin{proof}
We first prove that translated lifts have no poles along algebraic arcs starting in $U$. A compactification of the action graph will then turn this statement about orders into a locally uniform bound. Let $s\mapsto x(s)$ tend to $x_0\in U$, let $l(s)$ be a nonzero holomorphic lift with nonzero value at $s=0$, and let $t(s)\in T$ be meromorphic at zero. Since a torus is a product of multiplicative groups, write
\[
t(s)=\lambda(s)u(s)
\]
with an integral cocharacter $\lambda$ and a holomorphic map $u$ into $T$ whose value at zero is invertible. Replace $x(s),l(s)$ by $u(s)x(s),u(s)l(s)$. Rename their new initial point $u(0)x_0$ as $x_0$. It is still in $U$, which is invariant, so the weak-weight inequality remains available. It suffices to consider $t(s)=\lambda(s)$.

To compare the original lift with $\lambda(s)l(s)$, vary the rate at which its cocharacter parameter tends to zero. For rational $r\in[0,1]$, let $b(r)$ be the order of $\lambda(s^r)l(s)$ relative to a nonvanishing local frame at the limiting base point. Positive order means a zero of the lift and negative order means a pole. Take a common ramified parameter to define this order and divide by its degree. Projectivity gives the limiting base point, and algebraicity gives a finite integral order before normalization. Smoothness and positivity of the metric imply
\begin{equation}\label{eq:arcorder}
b(r)=\lim_{s\downarrow0}
\frac{-\log|\lambda(s^r)l(s)|_h^2}{-\log|s|^2}.
\end{equation}
For each fixed real $s>0$, the numerator is convex as a function of $r$. Indeed the pullback metric along the holomorphic orbit map from $\CC^*$ has psh weight, and invariance under the compact circle makes that weight convex in the real logarithmic coordinate. Composition with $r\mapsto r\log s$ preserves convexity. The positive denominator in \eqref{eq:arcorder} preserves it too. Passing to limits shows that $b$ is convex on rational points, and $b(0)=0$.

We claim that for sufficiently small positive rational $r$,
\begin{equation}\label{eq:initialslope}
b(r)=r w_L(x_0,\lambda).
\end{equation}
To prove the claim without assuming $L$ ample, first use the very ample bundles $J$ and $J+L$ from Section~\ref{subsec:weights}. Choose a nonvanishing local lift to $J$ along the arc and tensor it with $l(s)$ to obtain the lift to $J+L$. For either very ample bundle $C$, choose an eigenbasis of sections, with weights $\alpha_j$. Let $\nu_j\ge0$ be their orders relative to the initial lift; set $\nu_j=+\infty$ for an identically vanishing section. Relative to the transformed lift their orders are
\[
\nu_j-r\alpha_j.
\]
Because the sections generate $C$, the order of the transformed lift is
\[
-\min_j(\nu_j-r\alpha_j).
\]
For small enough $r>0$, every minimizer has $\nu_j=0$. The order is therefore $r\max_{\nu_j=0}\alpha_j=r w_C(x_0,\lambda)$. Subtract the formula for $J$ from that for $J+L$ to obtain \eqref{eq:initialslope}.

For such an $r$, convexity and $b(0)=0$ give
\[
b(r)\le r b(1),\qquad
b(1)\ge\frac{b(r)}r=w_L(x_0,\lambda)\ge0,
\]
where the last inequality is \eqref{eq:weakweights}. Thus a transformed lift has no pole along any such algebraic arc.

We now prove uniformity in $x$. Suppose \eqref{eq:orbitbound} were unbounded as $x\to x_0$. Compactify $T$ projectively, and compactify the total space of $L$ as $\mathbb P(\OO_Z\oplus L)$ with its infinity divisor. Consider the graph of the algebraic map $(t,x)\mapsto t\cdot l(x)$, retaining $x$ as a separate parameter, and take its closure in the resulting proper space over the parameter neighborhood. Since $Z$ is projective and $h$ is smooth, a sequence of unbounded norms gives a point of this closure at fiber infinity above $x_0$. Algebraic curve selection through that point, with generic point in the graph, gives after normalization an arc $x(s)\to x_0$ and a meromorphic torus arc $t(s)$ for which $t(s)l(x(s))$ has a pole. This contradicts the preceding paragraph. Hence the bound is locally uniform.

A holomorphic frame differs from an algebraic frame by a nowhere vanishing holomorphic function, bounded above and below on a smaller neighborhood. Positive powers preserve the estimate. This proves the remaining assertions.
\end{proof}

\subsection{The orbit-supremum metric}\label{subsec:minimum}

Fix a divisible positive integer $N$ for which \eqref{eq:canonicalquotient} is an isomorphism of line bundles. For $y\in W$ and a vector $v\in (NP)_y$, pull $v$ back to the orbit $q^{-1}(y)$ and define
\begin{equation}\label{eq:supnorm}
|v|_{h_{NP}}=\sup_{x\in q^{-1}(y)}|q^*v(x)|_{h^N}.
\end{equation}
The definition is independent of a representative of the orbit, since the isomorphism is equivariant. It scales by $|c|$ on replacing $v$ by $cv$. Finite stabilizer characters have been killed by the choice of $N$.

\begin{proposition}\label{prop:quotientmetric}
The norm \eqref{eq:supnorm} defines a metric whose $N$th root on $P$ has locally bounded psh weights on $W$, and whose pullback to every smooth resolution also has locally bounded psh weights. In particular $P$ is nef.
\end{proposition}

\begin{proof}
Consider a finite slice chart as in Proposition~\ref{prop:boundary}. Nearby quotient points have representatives on the slice near a fixed representative. Lemma~\ref{lem:orbitbound} gives a local upper bound for the norm of the pullback of a frame of $NP$. The supremum also includes the norm at the representative itself. The latter is bounded away from zero in a small chart, since the original metric is smooth and the frame is nonvanishing. Thus the supremum is finite and bounded above and below in frames. Its negative logarithm is consequently locally bounded.

To prove psh positivity, work over a small smooth slice coordinate ball $B$ and the universal cover $\CC^t\to T$. Let $e$ be the pulled-back local frame of $NP$ and set
\[
\Phi(z,\zeta)=-\log\bigl|\exp(\zeta)\cdot e(z)\bigr|_{h^N}^2
\qquad(z\in B,\ \zeta\in\CC^t).
\]
The action map is holomorphic and the translated frame is holomorphic and nonvanishing, so $\Phi$ is jointly psh. The metric is invariant under the compact form of $T$, hence $\Phi$ is independent of $\operatorname{Im}\zeta$. The domain $B\times\CC^t$ is pseudoconvex and has connected tube fibers. Kiselman's minimum principle \cite{Kiselman1978}, in the form of \cite[Theorem~1.1]{DWZZ2018}, therefore makes
\[
\inf_{\zeta\in\CC^t}\Phi(z,\zeta)
\]
psh or identically $-\infty$. It is the weight of \eqref{eq:supnorm}, and the bounds already proved exclude the latter possibility. Notice that the minimum principle supplies positivity after boundedness has been established; it does not supply the orbit bound.

The resulting slice weights are invariant under the finite slice group and have the correct transformations for the pullback of $NP$. They descend to locally bounded upper-semicontinuous functions on the finite quotient. To check positivity, take any holomorphic disc in that quotient. Normalize a component of its finite pullback to the slice that dominates the disc. After a finite ramified disc cover it has a holomorphic lift to the slice. The pulled-back weight is subharmonic there, hence subharmonic downstairs off the branch points; boundedness extends it across those points. This proves the disc test for the descended weight. In particular, after composition with any holomorphic disc in a smooth resolution, the same argument gives a subharmonic function. The pulled-back weights on the resolution are therefore psh. Equivalently, extend from the regular locus using the local upper bounds and \cite[Chapter~I, Theorem~5.24]{Demailly2012}, which is applied on the smooth resolution. Bounds above and below persist. Taking the $N$th root gives the claimed metric on $P$.

Finally let $C$ be an irreducible curve on a resolution. A local bounded psh weight restricts to a psh function on the normalization of $C$, rather than being identically $-\infty$. The resulting positive curvature measure has total mass equal to the degree of the restricted line bundle. This degree is nonnegative. Thus the pullback of $P$ has nonnegative intersection with every curve, so is nef on the projective resolution. Curves downstairs are dominated by curves upstairs; the projection formula gives nefness of $P$ on $W$.
\end{proof}

\subsection{The Euler premise and extension over the complement}\label{subsec:extension}

\begin{proof}[Proof of Theorem~\ref{thm:equivariant}]
Now assume that $Z$ is rationally connected. If the effective image of $T$ is trivial, every section is naturally invariant. Theorem~\ref{thm:klt} applied to $(Z,0)$ gives such a nonzero section, since $L$ is nef and has a smooth semipositive metric and $\chi(Z,\OO_Z)=1$.

Otherwise take the projective quotient $(W,D)$ of Proposition~\ref{prop:boundary}. It is $\QQ$-factorial and klt with effective rational boundary. Proposition~\ref{prop:quotientmetric} proves the nefness and bounded-metric hypotheses for $P=-(K_W+D)$.

Let $\widetilde W\to W$ be a smooth projective resolution. The rational map $Z\dashrightarrow\widetilde W$ is dominant. Resolve its indeterminacy to a morphism $\widetilde Z\to\widetilde W$. Rational connectedness is preserved by smooth projective birational modification and by a dominant rational map to a smooth projective target; see \cite{Kollar1996}. Alternatively, pullback of differential forms under this dominant morphism is injective and $H^0(\widetilde Z,\Omega^p_{\widetilde Z})=0$ for $p>0$. Hence $H^0(\widetilde W,\Omega^p_{\widetilde W})=0$ for $p>0$, and Hodge symmetry yields
\[
\chi(\widetilde W,\OO_{\widetilde W})=1.
\]
If $W$ is a point, this identity and the same conclusion hold directly.

Theorem~\ref{thm:klt} provides a nonzero section $\sigma$ of a positive Cartier multiple $mP$. Take a positive power so that $m$ is also divisible by the integer used in \eqref{eq:canonicalquotient}. Its pullback is a nonzero invariant section $\widetilde\sigma$ of $mL$ on $U$. The bounded quotient metric makes $|\sigma|$ bounded on the compact $W$: in finitely many relatively compact charts this follows from bounded holomorphic coefficients and the bounded frame norms. Definition~\eqref{eq:supnorm} implies
\begin{equation}\label{eq:extensionbound}
|\widetilde\sigma(x)|_{h^m}
\le |\sigma(q(x))|_{h_{mP}}\le C
\quad (x\in U).
\end{equation}

Let $x$ be any point of $Z\setminus U$ and choose a holomorphic nonvanishing frame of $mL$ on a small coordinate neighborhood. Smoothness of $h$ gives upper and lower positive bounds for the norm of that frame. The coefficient of $\widetilde\sigma$ is therefore a bounded holomorphic function off the algebraic set $Z\setminus U$. The removable-singularity theorem extends it holomorphically across that set \cite[Chapter~I, Corollary~5.25]{Demailly2012}. In particular the estimate applies at the generic point of every omitted divisor; no exception for an unstable divisor is made. The same argument, or normality, extends across the remaining codimension-two locus. Unique local extensions glue to a global section of $mL$.

For any $t\in T$, the translated global section agrees with the original on the dense open set $U$, so it agrees everywhere. The section remains nonzero since it was nonzero on $U$. This proves the invariant nonvanishing assertion.
\end{proof}

\begin{remark}\label{rem:character}
Natural linearization cannot be dropped even on $\mathbb P^1$. For the action $t\cdot[z_0:z_1]=[tz_0:z_1]$, the weights on $H^0(\mathbb P^1,-mK_{\mathbb P^1})$ under the natural linearization run from $-m$ to $m$. Tensor the linearization on $-K_{\mathbb P^1}$ by the character $t\mapsto t^2$. The weights in its $m$th power then run from $m$ to $3m$, so every positive power has zero invariant subspace. This is why the auxiliary character used to choose the stable open was not transferred to $L$.
\end{remark}
